% SOURCE_AUTHORITY: sga5_fr_workpass.tex, lines 14052--14080 (LF-normalized unit).
% SOURCE_SHA256: 791F4EFFC5E02832D5D77ED03518C8156D6F07E4C8238B03545DB93D883FBB28
\subsection*{Lema 6.4.5}
\emph{Se forem satisfeitas as relações seguintes:}
\begin{align*}
\nu_{x'}(f') &> 1 \text{,} \\
\nu_{x'}(g^{-1}f') &> 1 \quad ,
\end{align*}
\emph{então $g$ pertence ao primeiro grupo de ramificação $G_1$ (cf. 4.7):}
\[
g \in G_1
\]
\emph{isto é, $g$ induz a identidade sobre o espaço vetorial $\fm_{x'}/\fm_{x'}^2$.}

\subsection*{Demonstração}
Seja $\widehat{\mathcal{O}}_{x'} = k[[T]]$. O automorfismo de $\fm/\fm^2$ induzido por $g^{-1}$ é determinado por uma raiz da unidade $\varepsilon(g^{-1}) \in k$ pela condição
\[
g^{-1}T = \varepsilon(g^{-1})T \mod \fm^2 \quad ,
\]
onde $\fm$ é o ideal maximal do anel $k[[T]]$.

Por outro lado, a hipótese $\nu_{x'}(f') > 1$ dá
\[
f'(T) = T + a_\nu T^\nu + a_{\nu+1} T^{\nu+1} + \ldots \text{, } \quad \nu \geq 2 \quad .
\]

Por conseguinte,
\[
g^{-1}f^{\prime}(T) = \varepsilon(g^{-1})T + b_\nu T^\nu + b_{\nu+1} T^{\nu+1} + \ldots \text{,}
\]
e a hipótese $\nu(g^{-1}f^{\prime}) > 1$ implica $\varepsilon(g^{-1}) = 1$, c.q.d.
